\documentclass[10pt,a4paper]{amsart}
\usepackage[margin=19mm]{geometry}
\usepackage{amsmath,amssymb,mathtools}
\usepackage[scr=rsfs,cal=euler]{mathalfa}
\usepackage{xfrac}
\usepackage[unicode,hidelinks,bookmarks=false]{hyperref}
\newcommand{\ie}{i.e.\ }
\newcommand{\cf}{cf.\ }
\newcommand{\resp}{respectively\ }
\newcommand{\Subsection}{\subsection}
\providecommand{\nobreakdash}{\nobreak\hbox{-}}
\setlength{\emergencystretch}{2em}
\multlinegap0pt
\usepackage{bm}
\usepackage{bbm}
\usepackage{dsfont}
\usepackage{xspace}
\usepackage{cancel}
\usepackage{mathtools}
\usepackage{amsthm}
\makeatletter
\@ifundefined{Theorem}{\newtheorem{Theorem}{Theorem}}{}
\@ifundefined{Lemma}{\newtheorem{Lemma}[Theorem]{Lemma}}{}
\makeatother
\title[Isomonodromy and Painlev\'e Type Equations, Case Studies]{Isomonodromy and Painlev\'e Type Equations, Case Studies}
\author{Marius van der Put, Jaap Top}
\date{Source: arXiv:2404.15767v2 (CC BY 4.0)}
\begin{document}
\maketitle

\noindent\emph{Source and licence.} Isomonodromy and Painlev\'e Type Equations, Case Studies. Version arXiv:2404.15767v2 (CC BY 4.0), distributed under the Creative Commons Attribution 4.0 International licence, as recorded in the arXiv licence metadata for this exact version and shown on the abstract page https://arxiv.org/abs/2404.15767v2. Full rights notice: licenses/arxiv-2404.15767-CC-BY-4.0.txt.\par\smallskip

\noindent\textit{Editorial scope.} Counted unit: the Lemma whose source label is \texttt{lemma01} in the article (source0.tex lines 283--286, source label \texttt{lemma01}), delivered together with the article's own complete original proof of it (source0.tex lines 287--302, one \texttt{proof} environment of 16 lines). No mathematics is written, repaired or re-derived: statement and proof are the article's own text line for line. Required context retained uncounted: none. Every definition, display and label of those lines is retained unchanged. Omitted from this package and not redistributed: 1--282, 303--1736. Nothing inside a retained excerpt is omitted or altered: every retained body is delivered exactly as the article's lines are written, so no omission receipt of any kind accompanies this package. Citations: the retained excerpts carry no citation command at all, so no bibliography entry of the article is transcribed and no reference is redistributed. Reference adaptation: none was needed and none was made; every reference inside the delivered unit resolves inside it, so no unresolved reference or citation is produced. Printed slips kept literal: no printed word, symbol, hypothesis or number of the counted statement or of its proof was altered. Layout and class: the article's own class is \texttt{sigma} with packages including ifpdf; the wrapper is the lane's pinned \texttt{amsart} editorial wrapper at 10pt on A4, which restates the theorem-like environments the retained text needs and adds the article's own macro definitions that the retained text needs (none), with extra wrapper packages (bm, bbm, dsfont, xspace, cancel, mathtools). The delivered numbering is the unit's own. Assets: no retained span contains a figure environment, a graphics-inclusion command, a PDF-page inclusion, a table or a TikZ picture; no image file is copied into this package, the package declares no asset dependency and no image file is redistributed with it. Licence: version arXiv:2404.15767v2 of this article is distributed under the Creative Commons Attribution 4.0 International licence, as recorded in the arXiv licence metadata for this exact version and shown on the abstract page https://arxiv.org/abs/2404.15767v2. A full rights notice is delivered at licenses/arxiv-2404.15767-CC-BY-4.0.txt. Characters: every retained excerpt is pure ASCII, so the delivered engine, XeLaTeX with its default Latin Modern text font, reports no missing character.
\begin{Lemma} \label{lemma01} The monodromy space $\mathcal{R}$ is isomorphic to the quotient of the space $(\mathbb{C}^*)^{d-1}\times \mathbb{C}^N$ by the
 action of a~group isomorphic to $(\mathbb{C}^*)^{d-1}$. This quotient has an open, affine, dense subspace
 isomorphic to $(\mathbb{C}^*)^{d-1}\times \mathbb{C}^{N-d+1}$. In particular, $\dim \mathcal{R}=N$.
\end{Lemma}

\begin{proof} There is no restriction on the possibilities for the~\smash{${\rm St}_{d_j}$}. This produces the vector space~$\mathbb{C}^N$.
In order to make the restrictions on $\gamma$ explicit, we consider a Galois orbit, here written as,
$\overline{Q}=\{q_0,\dots ,q_{\ell -1}\}$ with $\dim V_{q_0}=f$. Choose a basis
$e_{1},\dots, e_{f}$ of eigenvectors for the action of $\gamma^\ell$ on $V_{q_0}$. Let $\alpha_1,\dots ,\alpha_f$ denote the
distinct eigenvalues. This basis is unique up to permuting and scaling of the basis vectors. Consider, for $i=1,\dots, \ell -1$, the basis of $V_{q_i}$ to be $\gamma ^i(e_1),\dots , \gamma ^i(e_f)$. One concludes that the data for the matrix of $\gamma$
on the space $\oplus V_{q_i}$ is equivalent to the tuple $(\alpha_1,\dots ,\alpha_f)$. Moreover, the set of all $\ell$th roots of all $\alpha _i$ is the set of eigenvalues of $\gamma$.

Thus the total data for $\gamma$ on $V$ is given by the eigenvalues of $\gamma^{l_i}$ on the space $ V_{q_{i,0}}$ for $i=1,\dots ,r$.
 Since, by assumption, $\gamma$ has determinant 1, the space of possibilities for $\gamma$ is~\smash{$(\mathbb{C}^*) ^{-1+\sum d_i}$}.


The automorphisms of $(V,\{V_q\},\gamma)$ are the $\tau \in {\rm PGL}(V)$ such that $\tau(V_q)=V_q$ for all $q$ and~${\tau \gamma=\gamma \tau}$. One concludes that $\tau$ is determined by its action on all $V_{q_{i,0}}$. Furthermore,
$\tau$~has on this space the same eigenvectors as $\gamma ^{\ell_i}$. This implies that the group of automorphism
is isomorphic to $(\mathbb{C}^*)^{d-1}$. It is seen that the group acts faithfully on the Stokes data. Finally, by scaling
suitable Stokes data to~1, one obtains this affine, open, dense subspace of $\mathcal{R}$.
\end{proof}

\end{document}
